% Pista geo-24s-q6 · Geometria
% Procedència: PAU Catalunya 2024 · Convocatòria de setembre · Sèrie 3 · Qüestió 6
\documentclass[11pt,a4paper]{article}
\usepackage{pista}
\pistainfo{Catalunya 2024}{Convocatòria de setembre}{Sèrie 3}

\begin{document}

\section*{\textcolor{blauPista}{Qüestió 6}}

\noindent Considereu el punt $P = (1, 3, 0)$ i el pla $\pi$ d'equació $x + 2y - 2z = -7$.

\subsection*{\textit{a)} \normalfont Sigui $r$ la recta que és perpendicular a $\pi$ i passa per $P$. Calculeu el punt d'intersecció de $\pi$ amb $r$. \hfill\textnormal{[1~punt]}}

\pista{Si $r$ és perpendicular a $\pi$, el seu vector director és, precisament, el vector normal de $\pi$ (que es llegeix directament dels coeficients de l'equació del pla). Amb el punt $P$ i aquest vector director, escriu $r$ en forma paramètrica. Per al punt d'intersecció, substitueix les expressions paramètriques de $r$ a l'equació del pla $\pi$, resol l'equació en el paràmetre, i d'aquesta manera trobaràs el punt concret.}

\subsection*{\textit{b)} \normalfont Calculeu la distància $d$ del punt $P$ al pla $\pi$. \hfill\textnormal{[0,5~punts]}}

\pista{Aplica directament la fórmula de la distància d'un punt a un pla.}

\subsection*{\textit{c)} \normalfont Calculeu l'equació d'un altre pla $\pi'$ que sigui paral·lel a $\pi$ i que també estigui a distància $d$ de $P$. \hfill\textnormal{[1~punt]}}

\pista{Pensa en la geometria de la situació: si $\pi'$ és paral·lel a $\pi$ i $P$ ha d'estar a la mateixa distància $d$ de tots dos, vol dir que $P$ queda just al mig entre $\pi$ i $\pi'$, un pla a cada costat de $P$. Com que $\pi'$ és paral·lel a $\pi$, té la mateixa forma $x + 2y - 2z + D' = 0$, on $D'$ és el terme independent que has de trobar. Imposa la condició que la distància de $P$ a $\pi'$ sigui $d$: obtindràs una equació amb un valor absolut, que tindrà dues solucions. Una d'elles correspon al pla $\pi$ original; l'altra és, precisament, la que et convé.}

\end{document}
